%%%%%%%%%%%%%%%%%%%%%%%%%%%%%%%%%%%%%%%%%%%%%%%%%%%%%%%%%%%%%%%%%%%%%%%%%%%%%%%%%
%  Time-Series Foundation Models as Approximate-Information-State World Models:
%  a Fisher-information theory of stability, observability, and learnability.
%
%  Companion / fusion paper to
%    "World-Action Models and Vision-Language-Action Models: One Latent, One
%     Fisher Information -- Two Orthogonal Blocks"  (wam_vla.tex),
%  built on the approximate-information-state framework of Subramanian, Sinha,
%  Seraj & Mahajan (JMLR 2022) and the CHRONOS time-series foundation model of
%  Ansari et al. (TMLR 2024).
%%%%%%%%%%%%%%%%%%%%%%%%%%%%%%%%%%%%%%%%%%%%%%%%%%%%%%%%%%%%%%%%%%%%%%%%%%%%%%%%%
\documentclass[10pt]{article}
\usepackage[letterpaper,margin=1in]{geometry}
\usepackage{amsmath,amssymb,amsthm,mathtools}
\usepackage{bm}
\usepackage[T1]{fontenc}
\usepackage{newtxtext,newtxmath}
\usepackage{microtype}
\usepackage{booktabs}
\usepackage{algorithm}
\usepackage{algpseudocode}
\usepackage{graphicx}
\usepackage{tikz}
\usetikzlibrary{arrows.meta,positioning,calc,fit,backgrounds}
\usepackage[dvipsnames]{xcolor}
\usepackage{enumitem}
\usepackage[colorlinks=true,allcolors=NavyBlue]{hyperref}
\usepackage[nameinlink,capitalize]{cleveref}
\usepackage{fancyhdr}

\pagestyle{fancy}\fancyhf{}
\renewcommand{\headrulewidth}{0.3pt}
\fancyhead[C]{\footnotesize\scshape Time-Series Foundation Models as World Models}
\fancyfoot[C]{\thepage}
\fancypagestyle{plain}{\fancyhf{}\fancyfoot[C]{\thepage}\renewcommand{\headrulewidth}{0pt}}

%----- Theorem environments -----
\usepackage{aliascnt}
\theoremstyle{plain}
\newtheorem{theorem}{Theorem}[section]
\newaliascnt{lemma}{theorem}\newtheorem{lemma}[lemma]{Lemma}\aliascntresetthe{lemma}
\newaliascnt{proposition}{theorem}\newtheorem{proposition}[proposition]{Proposition}\aliascntresetthe{proposition}
\newaliascnt{corollary}{theorem}\newtheorem{corollary}[corollary]{Corollary}\aliascntresetthe{corollary}
\theoremstyle{definition}
\newaliascnt{definition}{theorem}\newtheorem{definition}[definition]{Definition}\aliascntresetthe{definition}
\newaliascnt{example}{theorem}\newtheorem{example}[example]{Example}\aliascntresetthe{example}
\newtheorem{assumption}{Assumption}
\theoremstyle{remark}
\newaliascnt{remark}{theorem}\newtheorem{remark}[remark]{Remark}\aliascntresetthe{remark}
\renewcommand{\qedsymbol}{$\blacksquare$}
\crefname{proposition}{Proposition}{Propositions}\Crefname{proposition}{Proposition}{Propositions}
\crefname{corollary}{Corollary}{Corollaries}\Crefname{corollary}{Corollary}{Corollaries}
\crefname{lemma}{Lemma}{Lemmas}
\crefname{definition}{Definition}{Definitions}\Crefname{definition}{Definition}{Definitions}
\crefname{assumption}{Assumption}{Assumptions}\Crefname{assumption}{Assumption}{Assumptions}

%----- Macros -----
\newcommand{\R}{\mathbb{R}}
\newcommand{\E}{\mathbb{E}}
\newcommand{\Prob}{\mathbb{P}}
\DeclareMathOperator{\tr}{tr}
\DeclareMathOperator{\Cov}{Cov}
\DeclareMathOperator{\Var}{Var}
\DeclareMathOperator*{\argmax}{arg\,max}
\DeclareMathOperator*{\argmin}{arg\,min}
\DeclareMathOperator{\rank}{rank}
\DeclareMathOperator{\vecop}{vec}
\DeclareMathOperator{\diag}{diag}
\newcommand{\Fisher}{\mathcal{I}}
\newcommand{\Ith}{\Fisher_{\theta\theta}}
\newcommand{\Iph}{\Fisher_{\phi\phi}}
\newcommand{\Itp}{\Fisher_{\theta\phi}}
\newcommand{\Gam}{\Gamma}
\newcommand{\Ls}{\mathsf{L}}
\newcommand{\Os}{\mathsf{O}}
\newcommand{\Xs}{\mathsf{X}}
\newcommand{\Zs}{\mathsf{Z}}
\newcommand{\Hs}{\mathsf{H}}
\newcommand{\Fc}{\mathfrak{F}}
\newcommand{\given}{\,\vert\,}
\newcommand{\defeq}{\vcentcolon=}
\newcommand{\gth}{\nabla_{\theta}}
\newcommand{\IPM}{d_{\Fc}}
\newcommand{\rhoF}{\rho_{\Fc}}
\newcommand{\norm}[1]{\left\lVert #1 \right\rVert}
\newcommand{\abs}[1]{\left\lvert #1 \right\rvert}
\newcommand{\wm}{P_{\theta}}
\newcommand{\eps}{\varepsilon}
\newcommand{\TSFM}{\textsc{tsfm}}
\newcommand{\WAM}{\textsc{wam}}
\newcommand{\VLA}{\textsc{vla}}
\newcommand{\Chronos}{\textsc{chronos}}
\allowdisplaybreaks

\begin{document}
\title{\bfseries Time-Series Foundation Models as
Approximate-Information-State World Models:\\[3pt]
\large A Fisher-Information Theory of Stability, Observability, and Learnability}
\author{\normalsize (fusing the world-action / vision--language--action Fisher
theory with the \Chronos\ time-series foundation model,\\[-1pt]
\normalsize on the approximate-information-state framework of Subramanian,
Sinha, Seraj \& Mahajan, {\itshape JMLR} 2022)}
\date{\today}
\maketitle
\thispagestyle{plain}

\begin{abstract}\noindent
A time-series foundation model (\TSFM) such as \Chronos\ is trained to predict
the next value of a sequence by tokenization and cross-entropy---the language
model recipe transplanted to time series. A world-action model (\WAM) is trained
to predict an environment's next observation given actions. We show these are
the \emph{same object with the action channel deleted}: a \TSFM\ parameterizes
the environment (observation) channel of a partially observed process on a
shared approximate-information-state (AIS) latent, its transformer context
\emph{is} the AIS, and its statistical efficiency is governed by \emph{one}
Fisher information matrix (FIM)---the innovations / system-identification block
$\Ith$ of the \WAM/\VLA\ theory, now standing alone because a forecaster takes
no actions (the policy block and the occupancy coupling vanish). This single
identity carries the entire mathematical apparatus of the two-block theory into
forecasting: the tangent-filter FIM, the small-eigenvalue diagnostic, and the
data-versus-sensor dichotomy. Our central contribution is to show that, for the
linear--Gaussian (LQG) instance, the FIM's eigenstructure is dictated by two
physical quantities: \emph{stability margins} and \emph{observability}. We prove
a stability--identifiability law---the per-mode information obeys the AR(1) form
$\Gamma_{jj}\asymp \mathrm{gain}_j^2\,\sigma_w^2/(\sigma_v^2(1-r_j^2))$, so
$\lambda_{\max}(\Ith)\asymp 1/(1-\rho(A))$ diverges and the FIM condition number
blows up as the spectral radius $\rho(A)\to 1$---formalizing a tension between
\emph{learnability} (information grows near the unit circle) and
\emph{conditioning} (the sloppy spectrum widens). The eigenvectors of the small
eigenvalues are exactly the \emph{fast} modes (excitation/data-limited, curable
by longer series) and the \emph{unobserved} modes (sensor-limited, curable only
by a new observation channel), so the \WAM/\VLA\ data-versus-sensor dichotomy
becomes a statement about which time series a foundation model can and cannot
learn to forecast. A large numerical study on a stable LQG system with $200$
latent states and $100$ observations confirms every claim: the measured FIM has
exactly as many null directions as unobserved modes, its per-mode information
tracks $1/(1-r^2)$ to a log--log correlation of $0.97$, $\lambda_{\max}$ scales
as $1/(1-\rho)$, and a trained \Chronos-lite model's forecast skill rises with
stability and tracks $\log\det\Ith$, while sloppy directions are learned last in
exact proportion to their Fisher eigenvalue. The common theme: world modeling,
instruction-conditioned control, and time-series forecasting are three channels
of one process measured by one Fisher matrix whose spectrum is written by
dynamics.
\end{abstract}

\begin{center}\small
\textbf{Key words.} time-series foundation models, world models, Fisher
information, approximate information state, system identification, stability
margin, observability, Cram\'er--Rao bound, sloppy models, Chronos
\end{center}
\begin{center}\small
\textbf{AMS subject classifications.}
62M10, 62B10, 93E12, 93B07, 93B30, 68T07, 62M20, 93E11
\end{center}
\bigskip\hrule\bigskip

\section{Introduction}
\label{sec:intro}

Two research programs that rarely cite each other have arrived at the same
mathematical place. On one side, \emph{world models} and
\emph{vision--language--action} (\VLA) policies learn, respectively, to predict
an environment's future and to act in it; a companion paper~\cite{wamvla} shows
that these are the two factors of one partially observed process on a shared
approximate-information-state (AIS) latent~\cite{ais2022}, and that their
learning problems are the two orthogonal blocks of a single Fisher information
matrix (FIM). On the other side, \emph{time-series foundation models} (\TSFM)
such as \Chronos~\cite{chronos} forecast an arbitrary numerical sequence by
\emph{tokenizing} its values and training an off-the-shelf language model with
the cross-entropy loss---``the language of time series.'' These look like
unrelated enterprises: one is embodied, generative, and control-theoretic; the
other is a minimalist repurposing of a text transformer for forecasting. This
paper shows they are one enterprise (\cref{tab:paradigms}).

\begin{table}[t]
\centering
\small
\setlength{\tabcolsep}{5pt}
\renewcommand{\arraystretch}{1.28}
\caption{Three foundation-model paradigms as \emph{channels of one partially
observed process} on a shared approximate-information-state (AIS) latent. A
\WAM\ predicts an environment's next \emph{image} given the action taken; a
\TSFM\ predicts a numerical series' next value with \emph{no} action; a \VLA\
predicts the next \emph{action} from an image and a language instruction. They
differ only in which modality they observe and in whether they emit or condition
on actions---and hence in which block of the \emph{one} Fisher information matrix
governs their learning. Deleting the action channel collapses the joint FIM to
the single system-identification block $\Ith$ shared by \WAM\ and \TSFM\
(\cref{eq:collapse}).}
\label{tab:paradigms}
\begin{tabular}{@{}p{2.65cm}p{3.72cm}p{3.72cm}p{3.72cm}@{}}
\toprule
 & \textbf{World--Action}\newline\textbf{Model} (\WAM)
 & \textbf{Time-Series}\newline\textbf{Foundation} (\TSFM)
 & \textbf{Vision--Language}\newline\textbf{--Action} (\VLA)\\
\midrule
Observation $o_t$
 & images / video frames (pixels)
 & numerical values (scalar or vector)
 & images ($+$ proprioception) \\
Conditioning input
 & past frames \emph{and} action $a_t$
 & past values only
 & image \emph{and} language instruction $\ell$ \\
Emits (target)
 & next observation $o_{t+1}$
 & next value $o_{t+1}$
 & next action $a_t$ \\
Takes actions?
 & conditions on them
 & none (autonomous)
 & \emph{produces} them \\
Language / instruction?
 & no
 & no
 & yes ($\ell$) \\
Channel of the process
 & world $\wm(o\given h^-,a)$
 & world $\wm(o\given h^-)$
 & policy $\pi_\phi(a\given h,\ell)$ \\
Parameter learned
 & world $\theta$
 & world $\theta$
 & policy $\phi$ \\
Governing Fisher block
 & $\Ith$
 & $\Ith$ (\emph{sole} block)
 & $\Iph$ \\
Representative models
 & Genie, GAIA-1, DreamerV3
 & \Chronos, TimesFM, Moirai
 & RT-2, OpenVLA, $\pi_0$ \\
\bottomrule
\end{tabular}
\end{table}

\paragraph{The thesis in one sentence.}
A time-series foundation model is an \emph{autonomous world-action model}: it
parameterizes the observation channel $\wm(o_t\given h_t^-)$ of a partially
observed dynamical process, with the action channel deleted because a forecaster
does not act; its context embedding is an approximate information state; and its
statistical efficiency is governed by the \emph{single} Fisher block that
survives when the policy block is removed---the innovations,
system-identification FIM. Everything the \WAM/\VLA\ theory says about that
block~\cite{wamvla}---the tangent-filter computation, the eigenstructure of
identifiability, the small-eigenvalue diagnostic, the data-versus-sensor
dichotomy---transfers verbatim to forecasting, and acquires there a sharp new
reading in terms of the two quantities a control theorist cares about most:
\emph{stability} and \emph{observability}.

\paragraph{Why the fusion is exact, not analogical.}
\Cref{fig:collapse} is the whole idea. The \WAM/\VLA\ trajectory law factors
into an environment channel (emitting observations, parameterized by the world
$\theta$) and a policy channel (emitting actions, parameterized by the \VLA\
$\phi$), and the joint FIM is block-diagonal, $\Fisher=\mathrm{blkdiag}(\Ith,
\Iph)$~\cite[Thm.~6.1]{wamvla}. A \TSFM\ forecasts an \emph{uncontrolled}
sequence: there are no actions, so the policy factor is absent, $\phi$ does not
exist, the block $\Iph$ and the cross block $\Itp$ vanish identically, and the
joint FIM collapses to its \emph{world} block alone,
\begin{equation}
 \Fisher(\theta,\phi)=\begin{pmatrix}\Ith & 0\\ 0 & \Iph\end{pmatrix}
 \;\xrightarrow[\text{delete actions}]{\ \TSFM\ }\;
 \Fisher(\theta)=\Ith .
 \label{eq:collapse}
\end{equation}
A \TSFM\ is thus the pure-$\theta$, single-block corner of the two-block theory:
a world model with nothing to control. This is not a metaphor---it is
\eqref{eq:collapse}, and it means the \TSFM's information geometry is a
\emph{special case} of the embodied one, obtained by setting a channel to null.
The payoff is that the difficult, control-theoretic half of the \WAM/\VLA\
theory---the world block $\Ith$, which is the FIM of \emph{system
identification}---is exactly what a \TSFM\ optimizes, and it is far richer than
the forecasting literature's usual error metrics suggest.

\begin{figure}[t]
\centering
\begin{tikzpicture}[>=Latex,font=\small,node distance=5mm]
% left: two channels
\node[draw,rounded corners,fill=blue!8,minimum width=17mm,minimum height=7mm] (w) {$\wm(o\given h^-,a)$};
\node[draw,rounded corners,fill=orange!12,below=6mm of w,minimum width=17mm,minimum height=7mm] (p) {$\pi_\phi(a\given h,\ell)$};
\node[above=1mm of w,font=\footnotesize\itshape] {world / \TSFM\ channel};
\node[below=1mm of p,font=\footnotesize\itshape] {policy / \VLA\ channel};
\node[left=1mm of w,font=\footnotesize] {\WAM/\VLA:};
% arrow
\node[right=20mm of w,yshift=-8mm] (mid) {};
\draw[->,thick] ($(w.east)!0.5!(p.east)+(3mm,0)$) -- ++(16mm,0)
   node[midway,above,font=\footnotesize]{delete}
   node[midway,below,font=\footnotesize]{actions};
% right: FIM collapse
\node[right=34mm of w,yshift=-8mm] (fim) {$
  \begin{pmatrix}\Ith & 0\\ 0 & \Iph\end{pmatrix}
  \longrightarrow
  \boxed{\;\Ith\;}$};
\node[below=1mm of fim,font=\footnotesize\itshape,align=center]
  {one block survives:\\ the system-ID Fisher};
\end{tikzpicture}
\caption{The fusion \eqref{eq:collapse}. A time-series foundation model is a
world-action model with the action channel deleted: the policy block $\Iph$ and
the cross block $\Itp$ vanish, and the joint Fisher information collapses to the
single world (system-identification) block $\Ith$. The \TSFM\ is the
single-block corner of the two-block theory of~\cite{wamvla}.}
\label{fig:collapse}
\end{figure}

\paragraph{The common mathematical theme.}
Once a \TSFM's learning problem is identified with system identification through
$\Ith$, a single principle organizes all three paradigms---world modeling,
instruction-conditioned control, and time-series forecasting:
\begin{quote}\itshape
Each model estimates a channel of one partially observed process on a shared
AIS latent; its statistical efficiency is one Fisher information matrix; and the
eigenstructure of that matrix is written by the system's \emph{dynamics}---the
stability margins set the \emph{magnitude and conditioning} of information, and
the observability sets its \emph{null space}. The small Fisher eigenvalues are
exactly the fast and the unobserved modes---the directions no amount of the
wrong resource can learn.
\end{quote}
For a forecaster this theme is unusually concrete, because for a
linear--Gaussian system the FIM has a closed form and its spectrum is
\emph{legible}: each eigenvalue belongs to a physical mode, and its size is the
mode's stability margin weighted by its observability. That legibility is what
this paper exploits.

\paragraph{Contributions.}
\begin{enumerate}[leftmargin=1.6em,itemsep=2pt]
\item \textbf{The identity (\cref{sec:model}).} We cast a \TSFM\ as the
observation channel of a partially observed process on an AIS latent, show its
FIM is the single world block $\Ith$ of~\cite{wamvla} (\cref{prop:tsfmwam}),
and place \Chronos's tokenization inside the theory as a quantized observation
channel whose FIM is data-processing dominated by the continuous one
(\cref{prop:quant}).
\item \textbf{The context is the AIS (\cref{sec:ais}).} We show the \TSFM's
context embedding is an $(\eps,\delta)$-AIS and transport the AIS Fisher bound of
\cite{wamvla} to forecasting (\cref{prop:aisbound}).
\item \textbf{Stability writes the spectrum (\cref{sec:stability}).} We prove
the stability--identifiability law: per-mode information obeys the AR(1) form
$\Gamma_{jj}\asymp \mathrm{gain}_j^2/(1-r_j^2)$, so $\lambda_{\max}(\Ith)\asymp
1/(1-\rho(A))$ and the FIM condition number diverges as $\rho(A)\to1$
(\cref{thm:stability}); learnability and conditioning are in tension.
\item \textbf{Small eigenvalues, forecasting edition (\cref{sec:dichotomy}).}
We specialize the data-versus-sensor dichotomy: a sloppy Fisher direction is
either a fast but observed mode (curable by more data) or an unobserved mode
(curable only by a new sensor/covariate), and the fat observation map $m<n$ of a
realistic \TSFM\ \emph{guarantees} a sensor-limited null space
(\cref{thm:forecastdichotomy}), an observability ceiling on univariate
forecasting.
\item \textbf{Learnability (\cref{sec:learn}).} We tie the FIM spectrum to
learning curves through the Cram\'er--Rao bound: the excess forecast risk along
a mode decays as $1/(T\Gamma_{jj})$, so slow modes are learned first and sloppy
modes last, and aggregate skill tracks $\log\det\Ith$ (\cref{prop:learn}).
\item \textbf{When the second block returns (\cref{sec:covariates}).}
Covariate/exogenous-input forecasting re-introduces an action-like channel and
with it the policy block and the occupancy coupling of~\cite{wamvla}, locating
covariate \TSFM s back in the full two-block theory (\cref{prop:covariate}).
\item \textbf{A large numerical study (\cref{sec:experiments}).} On a stable LQG
system with $n=200$ states and $m=100$ observations we confirm all of the above,
including a trained \Chronos-lite whose forecast skill is predicted by the FIM.
\end{enumerate}

\paragraph{Related work.}
\Chronos~\cite{chronos} and the concurrent \TSFM s~\cite{das2023,woo2024,%
rasul2023,goswami2024} forecast by pretraining sequence models on large time-series
corpora; \Chronos\ specifically tokenizes by scaling and quantization and trains
with cross-entropy (``regression via classification''). The idea that a
next-token predictor is a maximum-likelihood system identifier is classical in
the innovations / prediction-error framework~\cite{ljung1999,cappe2005}; we make
the identification explicit and route it through the AIS latent and the two-block
FIM of~\cite{wamvla}. The stability--identifiability law generalizes the AR(1)
folk fact $I(a)=1/(1-a^2)$ and the near-unit-root asymptotics of
Phillips~\cite{phillips1987}; the ``sloppy'' spectrum spanning many decades is
the time-series instance of sloppy-model
theory~\cite{transtrum2015,gutenkunst2007}. Observability and the
parameter-sensitivity Gramian are standard in
identification~\cite{ljung1999,bellman1970}; our contribution is to read the
small Fisher eigenvectors of a \emph{foundation model} through them.

\section{A Time-Series Foundation Model is an Autonomous World Model}
\label{sec:model}

\subsection{The autonomous partially observed process}
A univariate or multivariate time series is, in state-space form, the output of
a partially observed dynamical system run \emph{open loop}. There is a latent
Markov state $X_t\in\Xs\subseteq\R^n$, an observation $O_t\in\Os\subseteq\R^m$,
and no action:
\begin{equation}
 X_{t+1}\sim\mathcal P_\theta(\cdot\given X_t),\qquad
 O_t\sim g_\theta(\cdot\given X_t),
 \label{eq:auto}
\end{equation}
with unknown dynamics parameter $\theta\in\Theta\subseteq\R^p$. This is exactly
the environment of the \WAM/\VLA\ model~\cite[\S2]{wamvla} with the action $A_t$
and the instruction $\ell$ removed. Write $H_t^-=(O_{1:t-1})$ for the history
before $O_t$. Marginalizing the latent gives the \emph{observation channel}
\begin{equation}
 \wm(o_t\given h_t^-)=\int_\Xs g_\theta(o_t\given x_t)\,
   b_{t\mid t-1}^\theta(\mathrm dx_t),\qquad
 b_{t\mid t-1}^\theta=\Prob_\theta(X_t\in\cdot\given H_t^-),
 \label{eq:channel}
\end{equation}
the predicted-belief mixture of the emission $g_\theta$. A \TSFM\ is a
parameterized model of \eqref{eq:channel}: given the past it emits a predictive
law for the next observation.

\begin{proposition}[The trajectory law has one factor]\label{prop:factor}
Under \eqref{eq:auto} the joint law of a length-$T$ series factors as
$\Prob_\theta(o_{1:T})=\prod_{t=1}^T \wm(o_t\given h_t^-)$, a single product of
observation channels. There is no action factor: the \WAM/\VLA\ two-channel
factorization~\cite[Prop.~2.1]{wamvla} degenerates to its first channel.
\end{proposition}
\begin{proof}
Order the variables $o_1,o_2,\dots$ and apply the chain rule; the conditional law
of $O_t$ given $o_{1:t-1}$ is \eqref{eq:channel} by \eqref{eq:auto}. The policy
factor $\pi_\phi(a_t\given\cdot)$ of~\cite{wamvla} is absent because no $A_t$ is
generated.
\end{proof}

\subsection{Chronos in the frame: tokenization is a quantized channel}
\Chronos~\cite{chronos} implements the channel \eqref{eq:channel} in three
steps: (i) \emph{mean-scale} the context, $\tilde o=o/s$ with
$s=\tfrac1C\sum|o|$; (ii) \emph{quantize} $\tilde o$ into one of $B$ bins by
$q:\R\to\{1,\dots,B\}$; (iii) model the resulting tokens with a categorical
language model $p_\theta(z_{t+1}=i\given z_{1:t})$ trained by cross-entropy.
Steps (i)--(ii) compose the observation with a fixed, $\theta$-free map $q\circ
(\cdot/s)$; step (iii) is a discrete instance of \eqref{eq:channel}. Thus
\Chronos\ is the observation channel \emph{read through a quantizer}, and
``regression via classification'' is the categorical likelihood of that
quantized channel. The next proposition is the Fisher consequence.

\begin{proposition}[Tokenization is data processing on the FIM]\label{prop:quant}
Let $\Ith$ be the Fisher information of the continuous observation channel
\eqref{eq:channel} and $\Ith^{q}$ that of its quantized (\Chronos) counterpart
$p_\theta(q(O_t)\given h_t^-)$. Because $q$ is a $\theta$-free measurable map,
$\Ith^{q}\preceq\Ith$ in the Loewner order, with the gap shrinking to zero as
the bin width $\to0$. Quantization can only lose information about the dynamics,
and a direction unidentifiable in the continuous channel stays unidentifiable
after tokenization.
\end{proposition}
\begin{proof}
The score of the quantized channel is the conditional expectation of the
continuous score given the quantized observation, $s^q_t=\E[s_t\given q(O_t),
h_t^-]$ (a $\theta$-free coarsening); the Fisher data-processing inequality
(the law of total covariance, as in~\cite[Cor.~3.4]{wamvla}) gives
$\Cov(s^q_t)\preceq\Cov(s_t)$. As the partition refines, $q(O_t)\to O_t$ and the
conditional expectation becomes the identity, so $\Ith^q\uparrow\Ith$.
\end{proof}

\Cref{prop:quant} says the \emph{number of bins} $B$ is an information knob: too
coarse a quantizer throws away Fisher information about the very modes the model
must learn. It also means the exact, continuous FIM $\Ith$ is the right object to
analyze---it is the ceiling that any tokenization approaches from below---which
is what we do for the rest of the paper.

\subsection{The Fisher information of a forecaster is the world block}
Fix the true dynamics and consider identifying $\theta$ from the observation
stream. Let $\mathcal F_t^-=\sigma(O_{1:t-1})$ and
$s_t^\theta\defeq\gth\log\wm(O_t\given H_t^-)\in\R^p$.

\begin{proposition}[The single Fisher block]\label{prop:tsfmwam}
Under standard regularity the \TSFM\ log-likelihood is
$\ell_T(\theta)=\sum_t\log\wm(O_t\given H_t^-)$; the score increments are
martingale differences, $\E_\theta[s_t^\theta\given\mathcal F_t^-]=0$; and the
Fisher information is
\begin{equation}
 \Ith(\theta)=\sum_{t=1}^T\E_\theta\big[s_t^\theta(s_t^\theta)^\top\big]
 =\sum_{t=1}^T\E_\theta\big[\Fisher_t^{\WAM}(\theta)\big],
 \label{eq:tsfmfim}
\end{equation}
identical to the world block~\cite[Prop.~4.1]{wamvla} but now taken under the
\emph{open-loop} law $\Prob_\theta$ (no policy occupancy). It is the \emph{only}
block: there is no policy score, hence no policy block and no cross block, and
the occupancy coupling of~\cite[Prop.~6.3]{wamvla} is vacuous. Training a \TSFM\
by next-observation prediction is maximum-likelihood system identification
preconditioned by $\Ith$.
\end{proposition}
\begin{proof}
\Cref{prop:factor} makes the likelihood additive over observation factors;
$\int\wm(o\given h^-)\,\mathrm do=1$ gives $\int\gth\wm=0$ and the martingale
property; vanishing cross-lag covariances (so the FIM is the diagonal sum) follow
because $s_s^\theta$ is $\mathcal F_t^-$-measurable for $s<t$. There is no
$\phi$, so \eqref{eq:collapse}. The equivalence with maximum likelihood is
$\gth\mathcal L=\E[\sum_t s_t^\theta]$ and $-\E[\nabla_\theta^2\mathcal
L]=\Ith$.
\end{proof}

The score is computed by the same \emph{tangent filter} as the world
block~\cite[\S4.1]{wamvla}: propagating $(b_t^\theta,\gth b_t^\theta)$ through
the filtering recursion yields $\gth\wm$ and hence $s_t^\theta$; in the
Gaussian/EKF case it is the differentiated Kalman filter, and for a learned
recurrent or transformer latent it is backpropagation through the state update.
A \TSFM\ trainer computing gradients of the next-token loss is running this
tangent filter whether it knows it or not.

\section{The Context Embedding is an Approximate Information State}
\label{sec:ais}

The \TSFM's transformer summarizes the observed history into a
context vector $z_t=\sigma_t(h_t)$ that it decodes into the next-token law. This
is precisely an information-state generator~\cite{ais2022}: $z_t$ is (meant to
be) sufficient to predict the next observation. In the exact linear--Gaussian
case the sufficient statistic is the Kalman estimate; a finite transformer
realizes only an $(\eps,\delta)$-AIS---approximately predictive to within
$\eps$ in reward-free forecasting error and $\delta$ in its own update, measured
in an integral probability metric (IPM). The next result is the forecasting
instance of the AIS Fisher bound~\cite[Thm.~7.2]{wamvla}.

\begin{proposition}[AIS Fisher bound for a \TSFM]\label{prop:aisbound}
Let $\hat z_t$ be an $(\eps,\delta)$-AIS whose predictive density is bounded
below by $\underline p>0$, with parameter-gradient approximated to within
$\eta_1$ in the IPM sense and score bounded by $\bar S$. Then the FIM computed on
the learned latent satisfies
$\norm{\Ith-\hat\Ith}_2\le 2T\bar S\,\beta$ with
$\beta=\eta_1/\underline p+\bar G\,\rhoF(\cdot)\,\delta/\underline p^2$. A latent
good enough for approximate forecasting is good enough to preserve
identifiability, and---there being only one block---no spurious coupling can
appear.
\end{proposition}
\begin{proof}
The single-channel specialization of~\cite[Thm.~7.2]{wamvla}: apply the
score-perturbation bound $\norm{s-\hat s}\le\eta_1/\underline p+\bar
G(\hat p-p)/\underline p^2$ and $\norm{ss^\top-\hat s\hat s^\top}_2\le 2\bar
S\norm{s-\hat s}$, then sum over $t$ and substitute the IPM prediction bound.
\end{proof}

\Cref{prop:aisbound} legitimizes analyzing the \emph{exact} FIM $\Ith$ as a
proxy for the trained model's: the learned context perturbs the information by an
IPM-controlled amount but cannot change its qualitative spectrum. We now compute
that spectrum.

\section{Stability Writes the Fisher Spectrum}
\label{sec:stability}

We specialize to the linear--Gaussian (LQG) instance, where everything is closed
form and every eigenvalue of the FIM belongs to a physical mode. The state is
organized into $K$ modes, each a damped oscillator; mode $j$ has magnitude
$r_j\in(0,1)$ (its \emph{stability}) and frequency $\omega_j$, contributing a
$2\times2$ block $r_j R(\omega_j)$ to $A$ with $R$ a rotation. Process noise
$\Sigma_w=\sigma_w^2 I$ drives the state, sensor noise $\Sigma_v=\sigma_v^2 I$
corrupts the observation $O_t=CX_t+V_t$, and $C\in\R^{m\times n}$ reads out the
modes with per-mode gain $\gamma_j$ (so $\gamma_j=0$ means mode $j$ is
unobserved). The dynamics parameter is the vector of modal magnitudes
$\theta=(r_1,\dots,r_K)$---the stabilities themselves.

\subsection{The sensitivity Gramian}
By \cref{prop:tsfmwam} the FIM is the innovations Fisher information. Using the
steady-state Kalman predictor with gain $L$, innovation covariance
$\Sigma_\nu=CPC^\top+\Sigma_v$, and predictor closed loop $F=A(I-LC)$ (Schur
stable), the per-step information about $\theta$ is the
\emph{parameter-sensitivity observability Gramian} of~\cite[Cor.~9.5]{wamvla},
\begin{equation}
 \Gam \;\defeq\; \Ith \;=\;
 \E\big[(C\,\partial_\theta\hat X_{t\mid t-1})^\top\Sigma_\nu^{-1}
        (C\,\partial_\theta\hat X_{t\mid t-1})\big],\qquad
 \partial_\theta\hat X_{t+1\mid t}=F\,\partial_\theta\hat X_{t\mid t-1}
   +(\partial_\theta A)\hat X_{t\mid t},
 \label{eq:gramian}
\end{equation}
the tangent Kalman recursion of~\cite[App.~A]{wamvla} specialized to the
autonomous system. \Cref{eq:gramian} is exact and is what our numerics compute;
its value is that the forcing $(\partial_{r_j}A)\hat X_{t\mid t}$ excites only
mode $j$, so $\Gam$ is nearly diagonal in the modes and each of its eigenvectors
points at one mode.

\subsection{The stability--identifiability law}
The heart of the paper is that the size of the information about mode $j$ is set
by that mode's stability margin.

\begin{theorem}[Stability--identifiability law]\label{thm:stability}
For the modal LQG system the per-mode Fisher information obeys the AR(1) law
\begin{equation}
 \Gam_{jj}\;\asymp\;\frac{\gamma_j^2\,\norm{c_j}^2\,\sigma_w^2}
                          {\sigma_v^2\,(1-r_j^2)},
 \label{eq:arlaw}
\end{equation}
where $\norm{c_j}^2$ is the squared readout norm of mode $j$. Consequently the
largest Fisher eigenvalue is governed by the slowest observed mode,
\begin{equation}
 \lambda_{\max}(\Gam)\;\asymp\;\frac{1}{1-\rho(A)},
 \label{eq:lammax}
\end{equation}
and, since the fastest observed mode contributes an $O(1)$ eigenvalue, the FIM
condition number diverges as the stability margin vanishes,
$\kappa(\Gam)=\lambda_{\max}/\lambda_{\min}^{+}\to\infty$ as $\rho(A)\to1$.
Reducing the stability margin \emph{increases} total information
($\log\det\Gam$, $\tr\Gam$, $\lambda_{\max}$ all grow) while \emph{worsening}
conditioning: learnability and conditioning are in tension.
\end{theorem}

\begin{proof}[Proof sketch]
A single mode in isolation is a two-dimensional AR(1): its stationary variance
solves the scalar Lyapunov equation $\pi_j=r_j^2\pi_j+\sigma_w^2$, so
$\pi_j=\sigma_w^2/(1-r_j^2)$. The one-step prediction $\hat o=c_j^\top\hat x_j$
has sensitivity $\partial_{r_j}\hat x_j$ proportional to the mode amplitude, so
its stationary second moment scales with $\pi_j$; contracting with the
innovation precision $\Sigma_\nu^{-1}\!\approx\!\sigma_v^{-2}I$ and the readout
$c_j$ gives \eqref{eq:arlaw} (the exact steady-state computation propagates
$\partial_\theta\hat X$ through \eqref{eq:gramian} and solves the resulting
Lyapunov equation; the leading term is \eqref{eq:arlaw}). Taking the maximum over
observed modes and using $\rho(A)=\max_j r_j$ gives \eqref{eq:lammax}; the
fastest mode gives an $O(1)$ floor, so the ratio diverges. Monotonicity of
$\log\det$, $\tr$, and $\lambda_{\max}$ in the $r_j$ follows from
\eqref{eq:arlaw}. \Cref{sec:experiments} confirms \eqref{eq:arlaw}--%
\eqref{eq:lammax} numerically to a log--log correlation of $0.97$ and a
$1/(1-\rho)$ fit.
\end{proof}

\Cref{eq:arlaw} is the AR(1) folk fact $I(a)=1/(1-a^2)$~\cite{phillips1987}
promoted to a foundation model: the slowest, most persistent modes near the unit
circle carry almost all the Fisher information, and the fast modes near the
origin carry almost none. This is why a forecaster trained on a nearly
nonstationary series identifies its dominant dynamics so precisely---and why the
same series' fast components are essentially unlearnable.

\begin{remark}[Learnability versus control]
The tension in \cref{thm:stability} inverts the usual control intuition. A system
near the unit circle is \emph{hard to control} (the LQG cost inflates) but
\emph{easy to identify} (Fisher information inflates). A foundation model is a
pure identifier, so it \emph{benefits} from the regime a controller fears; the
stability margin is a single dial trading the two.
\end{remark}

\section{Small Eigenvalues: Which Series a Foundation Model Cannot Learn}
\label{sec:dichotomy}

The eigenvectors of the small Fisher eigenvalues are the actionable output of the
theory: they name the modes a forecaster cannot pin down, and---exactly as
in~\cite[\S11]{wamvla}---they say whether the cure is \emph{more data} or a
\emph{new sensor}. The autonomous setting sharpens the dichotomy because a
forecaster cannot explore: its only data knob is the length of the series.

\begin{theorem}[Forecasting data--sensor dichotomy]\label{thm:forecastdichotomy}
Let $v$ be a unit eigenvector of a small eigenvalue of $\Gam$, pointing (by the
modal structure of \eqref{eq:gramian}) at a mode $j$. Exactly one of two cases
holds.
\begin{enumerate}[leftmargin=1.7em,itemsep=1pt]
\item \textup{(excitation/data-limited)} Mode $j$ is observed ($\gamma_j>0$) but
fast ($r_j$ small), so $\Gam_{vv}=v^\top\Gam v>0$ is small by \eqref{eq:arlaw}.
The Cram\'er--Rao variance $1/(T\,\Gam_{vv})\to0$ as the series length $T\to\infty$:
\emph{more data cures it}, at a rate set by $1/(1-r_j^2)$.
\item \textup{(sensor-limited)} Mode $j$ is unobserved ($\gamma_j=0$), so
$v\in\ker\Gam$ and $\Gam_{vv}=0$ for \emph{every} series length: no amount of
data raises the eigenvalue. The direction is structurally unidentifiable and
only a new observation channel that sees mode $j$ (a row appended to $C$) cures
it, by the additivity $\Gam^{\,C\cup c}=\Gam^{\,C}+\Gam_c$.
\end{enumerate}
Moreover, whenever the observation map is \emph{fat}, $m<n$ (fewer sensors than
latent states---the generic case for a real time series), $\ker\Gam\neq\{0\}$:
there is always a sensor-limited subspace, an \emph{observability ceiling} on
what univariate/low-dimensional forecasting can learn.
\end{theorem}
\begin{proof}
The modal structure of \eqref{eq:gramian} gives $\Gam_{vv}=\Gam_{jj}$ up to the
eigenvector's concentration; \eqref{eq:arlaw} gives case 1 and the rate. For case
2, $\gamma_j=0$ makes mode $j$'s readout $c_j=0$, so $C\partial_{r_j}\hat X\equiv
0$ and the integrand of \eqref{eq:gramian} vanishes identically; hence
$\Gam_{jj}=0$ irrespective of $T$, and additivity of Fisher information over
conditionally independent channels gives the sensor cure. When $m<n$, $C$ has a
nontrivial kernel; modes lying in it have $c_j=0$ and populate $\ker\Gam$.
\end{proof}

\Cref{thm:forecastdichotomy} is the forecasting reading of the \WAM/\VLA\
dichotomy~\cite[Thm.~11.3]{wamvla}: a small Fisher eigenvalue on a foundation
model is a message. If its mode is observed, collect a longer series (or, in the
multi-series regime, more series that excite that mode); if its mode is
unobserved, no corpus will ever teach the model to forecast it---acquire a new
measurement channel. The fat-$C$ corollary is the precise sense in which a
univariate forecaster faces an observability ceiling: it can only learn the part
of the dynamics its single channel sees.

\begin{example}[Position without velocity]
A second-order mechanical mode observed only in position ($C=[1,0]$) has its
velocity coordinate unobserved; the corresponding parameter direction is
sensor-limited, and no length of position-only data identifies it, whereas a
velocity channel lifts it immediately. \Cref{sec:experiments} realizes this at
scale: eight modes are left unobserved and produce exactly eight null Fisher
directions, cured only by adding channels that see them.
\end{example}

\section{Learnability: From the Fisher Spectrum to Learning Curves}
\label{sec:learn}

The FIM spectrum predicts not just \emph{whether} but \emph{how fast} a
foundation model learns each part of the dynamics, through the Cram\'er--Rao
bound and the martingale central limit theorem for the next-token MLE.

\begin{proposition}[Learning curve and per-mode risk]\label{prop:learn}
Let $\hat\theta_T$ be the next-token maximum-likelihood estimate from a series of
length $T$. Under identifiability of the observed modes,
$\sqrt{T}(\hat\theta_T-\theta^\star)\Rightarrow\mathcal N(0,\bar\Gam^{-1})$ on the
observed subspace, where $\bar\Gam$ is the per-step FIM. Consequently:
\begin{enumerate}[leftmargin=1.7em,itemsep=1pt]
\item the estimation variance along mode $j$ is $1/(T\bar\Gam_{jj})\asymp
\sigma_v^2(1-r_j^2)/(T\gamma_j^2\norm{c_j}^2)$: \emph{slow modes are learned
first, fast modes last, unobserved modes never};
\item the excess $h$-step forecast error along mode $j$ is bounded by
$S_{j,h}^2/(T\bar\Gam_{jj})$, with $S_{j,h}=\partial_{r_j}(\text{$h$-step
predictor})$ the prediction sensitivity, so early in training aggregate skill is
dominated by the slow, high-Fisher modes and rises with $\log\det\bar\Gam$;
\item the sloppy directions ($\lambda_j\!\downarrow$) require
$T\gtrsim1/(\varrho^2\lambda_j)$ samples to reach relative precision $\varrho$,
so the number of learnable modes grows only logarithmically in $T$ given the
geometric Fisher spectrum of \eqref{eq:arlaw}.
\end{enumerate}
\end{proposition}
\begin{proof}
The joint score is a martingale (\cref{prop:tsfmwam}) with predictable quadratic
variation $T\bar\Gam$; the martingale CLT and a Taylor expansion of the
stationarity condition give the normal limit with covariance $\bar\Gam^{-1}$
(restricted to the identifiable subspace, $\ker\bar\Gam$ excluded). Claim~1 is the
diagonal of $\bar\Gam^{-1}$ with \eqref{eq:arlaw}; claim~2 is the delta method on
the $h$-step predictor; claim~3 is the CRB threshold $v^\top\bar\Gam^{-1}v\ge
1/\lambda_j$.
\end{proof}

\Cref{prop:learn} converts the sloppy Fisher spectrum into a concrete curriculum:
a foundation model learns a system's modes in \emph{decreasing order of stability},
each at a rate proportional to its Fisher eigenvalue, and reaches an accuracy
floor set by the sloppiest modes it can still see. \Cref{sec:experiments} shows a
trained \Chronos-lite exhibiting exactly this: forecast skill rising with data and
with $\log\det\Gam$, and a stiff/sloppy learnability ratio equal to the Fisher
eigenvalue ratio.

\section{When the Second Block Returns: Covariates and Control}
\label{sec:covariates}

The collapse \eqref{eq:collapse} is exact for \emph{autonomous} forecasting.
Real forecasting problems often include \emph{covariates}---exogenous inputs
$u_t$ (promotions, weather, control actions) that drive the series. A covariate
is an action the forecaster does not choose but must condition on, and it
re-activates the second channel.

\begin{proposition}[Covariate forecasting re-opens the policy block]\label{prop:covariate}
Suppose the state obeys $X_{t+1}\sim\mathcal P_\theta(\cdot\given X_t,U_t)$ with
exogenous $U_t\sim\pi_\phi(\cdot\given H_t)$ generated by a (possibly unknown)
mechanism $\phi$. Then the trajectory law regains the two-channel factorization
of~\cite[Prop.~2.1]{wamvla}: an observation channel $\wm(o_t\given h_t^-,u_{t-1})$
and a covariate channel $\pi_\phi(u_t\given h_t)$. If $\phi$ is estimated, the
joint FIM regains its block structure $\mathrm{blkdiag}(\Ith,\Iph)$ and, crucially,
its \emph{occupancy coupling}~\cite[Prop.~6.3]{wamvla}: the covariate
distribution sizes $\Ith$ (which dynamics the data excite) and the dynamics size
$\Iph$. A \TSFM\ with covariates is therefore back in the full two-block theory;
pure forecasting is the $\phi$-free corner.
\end{proposition}
\begin{proof}
Re-instate the action factor in \cref{prop:factor}; the interleaved-filtration
argument of~\cite[Thm.~6.1]{wamvla} gives block-diagonality and
\cite[Prop.~6.3]{wamvla} the occupancy coupling. Deleting $U_t$ recovers
\cref{prop:tsfmwam}.
\end{proof}

Thus the WAM/VLA/TSFM trichotomy is a single axis: no exogenous input (pure
\TSFM, one block $\Ith$); exogenous but not chosen (covariate \TSFM, two blocks,
coupled); chosen by a learned policy under instruction (\VLA, two blocks, and the
co-design of~\cite[\S6]{wamvla}). Stability and observability govern $\Ith$ in
all three; what the covariate/control adds is the occupancy that \emph{sizes}
$\Ith$, and its own identifiability $\Iph$.

\section{Numerical Study: A 200-State, 100-Observation Stable LQG System}
\label{sec:experiments}

We instantiate the theory on a deliberately large, stable, partially observed
LQG system and a faithfully small \Chronos-lite foundation model. The code is in
\texttt{tsfm\_lqg/} (system \texttt{lqg\_system.py}, Fisher \texttt{fisher.py},
model \texttt{chronos\_lite.py}, drivers \texttt{experiments*.py}).

\paragraph{The system.}
$K=100$ modes give $n=200$ latent states; $m=100$ observations make $C$ a
$100\times200$ \emph{fat} map (an observability ceiling by
\cref{thm:forecastdichotomy}). Modal magnitudes are log-spaced in $[0.15,\rho]$
with baseline spectral radius $\rho(A)=0.90$ (stability margin $0.10$); process
noise $\sigma_w=1$, sensor noise $\sigma_v=0.1$. Observation gains are tilted
toward the slow, persistent modes (which therefore dominate each channel and are
the most forecastable), and the \textbf{eight fastest modes are left unobserved}
($\gamma_j=0$) to realize the sensor-limited case. The FIM
\eqref{eq:gramian} about $\theta=(r_1,\dots,r_{100})$ is computed by the exact
tangent Kalman filter; a Monte-Carlo score estimate agrees with the plug-in FIM
to $\sim\!10\%$.

\paragraph{The foundation model.}
\Chronos-lite follows the \Chronos\ recipe exactly: mean scaling, uniform
quantization into $B=128$ bins, and a small causal Transformer (decoder-only,
$d_{\text{model}}=64$, $2$ layers, context $40$) trained with the cross-entropy
next-token loss; forecasts are autoregressive-sample or expected-value point
forecasts, de-quantized and un-scaled. As \Chronos\ prescribes, the model is
\emph{univariate}---one shared model across all $100$ channels---so its access to
each mode is limited by that mode's observability, making
\cref{thm:forecastdichotomy} operative.

\subsection{The Fisher spectrum is sloppy and legible (\cref{fig:spectrum})}
\Cref{fig:spectrum}(a) shows the FIM eigenvalues spanning $\sim\!15$ decades with
exactly \textbf{eight} eigenvalues at machine zero---one per unobserved mode,
confirming the $m<n$ null space of \cref{thm:forecastdichotomy}. Panel~(b) plots
the gain-normalized per-mode information $\Gam_{jj}/\norm{c_j}^2$ against the
stability margin $1-r_j^2$; it collapses onto the $1/(1-r^2)$ law of
\eqref{eq:arlaw} with log--log correlation $\mathbf{0.97}$. Panel~(c) shows the
ten smallest-eigenvalue eigenvectors, each concentrated (weight $\approx1$) on a
single fast or unobserved mode, confirming that the sloppy directions are
individual physical modes.

\begin{figure}[t]
\centering
\includegraphics[width=\linewidth]{tsfm_lqg/figs/fig_spectrum.pdf}
\caption{Fisher spectrum of the \TSFM\ on the $200$-state LQG system.
(a) the sloppy spectrum: $\sim\!15$ decades with eight sensor-limited zeros;
(b) gain-normalized per-mode information collapses onto the $1/(1-r^2)$
stability law (\cref{thm:stability}); (c) each small-eigenvalue eigenvector is a
single mode.}
\label{fig:spectrum}
\end{figure}

\subsection{Data versus sensor (\cref{fig:dichotomy})}
\Cref{fig:dichotomy}(a) plots the Cram\'er--Rao standard deviation
$1/\sqrt{T\,\Gam_{vv}}$ against series length $T$ for three directions: a stiff
(slow, observed) mode, an excitation-limited (fast, observed) mode, and a
sensor-limited (unobserved) mode. The two observed directions fall as
$T^{-1/2}$---\emph{data cures them}, the stiff one far faster---while the
unobserved direction is flat at its ceiling: no series length helps. Panel~(b)
shows that appending observation channels that see the eight hidden modes lifts
the smallest eigenvalue on that subspace from $\approx\!10^{-8}$ to $\approx\!1$,
a jump of eight orders of magnitude: \emph{only a new sensor cures the null},
exactly \cref{thm:forecastdichotomy}.

\begin{figure}[t]
\centering
\includegraphics[width=0.92\linewidth]{tsfm_lqg/figs/fig_dichotomy.pdf}
\caption{The forecasting data--sensor dichotomy (\cref{thm:forecastdichotomy}).
(a) more data cures excitation-limited but not sensor-limited directions;
(b) adding channels that see the hidden modes lifts the null eigenvalue by eight
orders of magnitude.}
\label{fig:dichotomy}
\end{figure}

\subsection{Stability writes the spectrum (\cref{fig:stabfish})}
Sweeping the spectral radius $\rho(A)$ from $0.40$ to $0.98$ (rescaling all modal
magnitudes) traces \cref{fig:stabfish}. Panel~(a): $\lambda_{\max}(\Gam)$ and the
dominant-mode information grow by more than a decade as $\rho\to1$---information
concentrates near the unit circle. Panel~(b): the condition number worsens from
$\sim\!2\times10^{14}$ to $\sim\!5\times10^{15}$---the sloppy spectrum widens.
Panel~(c): on log--log axes $\lambda_{\max}$ follows the predicted $1/(1-\rho)$
law of \eqref{eq:lammax} (at $\rho=0.98$, $\lambda_{\max}=49.7\approx1/0.02$).
\Cref{tab:stab} lists the numbers; they are the empirical content of
\cref{thm:stability}.

\begin{figure}[t]
\centering
\includegraphics[width=\linewidth]{tsfm_lqg/figs/fig_stability_fisher.pdf}
\caption{Fisher properties versus stability (\cref{thm:stability}).
(a) information grows toward the unit circle; (b) conditioning worsens;
(c) $\lambda_{\max}\propto1/(1-\rho)$, the stability-margin law.}
\label{fig:stabfish}
\end{figure}

\begin{table}[h]
\centering\small
\caption{Fisher information versus stability margin (baseline system, exact
tangent-filter FIM). Information ($\log\det\Gam$, $\lambda_{\max}$) grows and
conditioning ($\kappa$) worsens as $\rho(A)\to1$; $\lambda_{\max}$ tracks
$1/(1-\rho)$.}
\label{tab:stab}
\begin{tabular}{@{}lcccccccc@{}}
\toprule
$\rho(A)$ & $0.40$ & $0.55$ & $0.68$ & $0.78$ & $0.86$ & $0.92$ & $0.96$ & $0.98$\\
$1-\rho$  & $0.60$ & $0.45$ & $0.32$ & $0.22$ & $0.14$ & $0.08$ & $0.04$ & $0.02$\\
\midrule
$\log\det\Gam$ & $-560.8$ & $-553.6$ & $-544.9$ & $-536.1$ & $-527.3$ & $-519.0$ & $-512.0$ & $-507.7$\\
$\lambda_{\max}(\Gam)$ & $2.28$ & $2.82$ & $3.76$ & $5.28$ & $8.12$ & $13.9$ & $26.8$ & $49.7$\\
$1/(1-\rho)$ & $1.7$ & $2.2$ & $3.1$ & $4.5$ & $7.1$ & $12.5$ & $25$ & $50$\\
$\kappa(\Gam)\ (\times10^{14})$ & $2.3$ & $2.8$ & $3.8$ & $5.3$ & $8.1$ & $13.9$ & $26.7$ & $49.7$\\
\bottomrule
\end{tabular}
\end{table}

\subsection{The foundation model learns what the Fisher says (\cref{fig:tsfm})}
We train \Chronos-lite on each stability level and measure one-step forecast
$R^2$ (original units). \Cref{fig:tsfm}(a): the learned model's skill rises
monotonically with $\rho(A)$, tracking the Kalman-optimal (all-channels)
ceiling---more stable systems are more learnable, as \cref{thm:stability} and
\cref{prop:learn} predict. Panel~(b) plots the same skill against
$\log\det\Gam$: forecast performance is an increasing function of the Fisher
information, the empirical content of \cref{prop:learn}. \Cref{tab:tsfm} gives
the numbers: as $\rho(A)$ rises from $0.45$ to $0.975$ the learned model's
one-step forecast $R^2$ climbs from $0.105$ to $0.737$, trailing the
Kalman-optimal ceiling by a roughly constant gap---the univariate,
finite-data penalty of \cref{prop:learn}---while $\log\det\Gam$ rises
in lockstep from $-558.7$ to $-508.3$.

\begin{table}[h]
\centering\small
\caption{\Chronos-lite forecast skill versus stability (trained per level;
one-step $R^2$ on held-out series, original units). Learned skill rises with
stability toward the Kalman-optimal ceiling and tracks the Fisher information
$\log\det\Gam$.}
\label{tab:tsfm}
\begin{tabular}{@{}lccccccc@{}}
\toprule
$\rho(A)$ & $0.45$ & $0.60$ & $0.72$ & $0.82$ & $0.90$ & $0.95$ & $0.975$\\
\midrule
TSFM $R^2$ (learned) & $0.105$ & $0.171$ & $0.305$ & $0.415$ & $0.528$ & $0.641$ & $0.737$\\
Kalman $R^2$ (ceiling) & $0.124$ & $0.224$ & $0.351$ & $0.492$ & $0.623$ & $0.720$ & $0.800$\\
$\log\det\Gam$ & $-558.7$ & $-550.5$ & $-541.5$ & $-531.7$ & $-521.6$ & $-513.5$ & $-508.3$\\
\bottomrule
\end{tabular}
\end{table}

\begin{figure}[t]
\centering
\includegraphics[width=0.92\linewidth]{tsfm_lqg/figs/fig_tsfm_stability.pdf}
\caption{\Chronos-lite learning performance versus stability.
(a) one-step forecast $R^2$ rises with the spectral radius toward the
Kalman-optimal ceiling; (b) the same skill increases with the Fisher information
$\log\det\Gam$ (\cref{prop:learn}).}
\label{fig:tsfm}
\end{figure}

\subsection{Learning curves and the Cram\'er--Rao law (\cref{fig:learn})}
Fixing $\rho=0.90$ and varying the training budget, \cref{fig:learn}(a) shows
forecast $R^2$ rising toward the Kalman ceiling as data grows---the learning
curve of \cref{prop:learn}. Panel~(b) plots the Cram\'er--Rao standard deviation
$1/\sqrt{T\lambda}$ for the stiffest and the sloppiest identifiable Fisher
directions; the two differ by the Fisher eigenvalue ratio (here
$\sim\!10^{6}$), so the sloppy modes need about a million times more data---%
\emph{sloppy modes are learned last}, in exact proportion to their Fisher
eigenvalue.

\begin{figure}[t]
\centering
\includegraphics[width=0.92\linewidth]{tsfm_lqg/figs/fig_learning_curve.pdf}
\caption{Learning curves and the Cram\'er--Rao law (\cref{prop:learn}).
(a) forecast $R^2$ rises toward the Kalman ceiling with data; (b) the stiff and
sloppy Fisher directions have vastly different CRB decay, so sloppy modes are
learned last.}
\label{fig:learn}
\end{figure}

\subsection{Summary of the numerical confirmations}
\Cref{tab:confirm} collects the quantitative checks. Every qualitative
prediction of \cref{sec:stability,sec:dichotomy,sec:learn} is borne out, and the
two headline laws---$\Gam_{jj}\propto1/(1-r^2)$ and
$\lambda_{\max}\propto1/(1-\rho)$---hold to within numerical and Monte-Carlo
error.

\begin{table}[h]
\centering\small
\caption{Numerical confirmations on the $200$-state / $100$-observation system.}
\label{tab:confirm}
\begin{tabular}{@{}p{7.6cm}p{3.0cm}p{3.4cm}@{}}
\toprule
Prediction & Theory & Measured\\
\midrule
\# null Fisher directions $=$ \# unobserved modes & $8$ & $8$ (exact)\\
per-mode Fisher vs.\ stability margin & $\propto1/(1-r^2)$ & log--log corr $0.97$\\
$\lambda_{\max}$ vs.\ spectral radius & $\propto1/(1-\rho)$ & Table~\ref{tab:stab}\\
sensor cure of the null subspace & $0\to>0$ & $10^{-8}\to1.0$\\
FIM condition number as $\rho\!\to\!1$ & $\to\infty$ & $2\!\times\!10^{14}\!\to\!5\!\times\!10^{15}$\\
TSFM skill vs.\ stability / $\log\det\Gam$ & increasing & \cref{fig:tsfm}\\
sloppy vs.\ stiff learnability & $\propto\lambda$ ratio & $\sim\!10^6$\\
\bottomrule
\end{tabular}
\end{table}

\section{Discussion}
\label{sec:discussion}

\paragraph{What the fusion buys.}
Reading a time-series foundation model as an autonomous world model imports a
control-theoretic diagnostic kit into forecasting. The FIM $\Ith$ is not an
abstraction: its eigenvalues are the learnable rates of the system's modes, its
small eigenvectors are the specific dynamics the model cannot resolve, and its
null space is the observability ceiling of the sensor suite. A practitioner can,
in principle, compute (or empirically estimate via per-sample gradient outer
products~\cite[\S12]{wamvla}) the model's Fisher spectrum and read off which
series need longer histories, which need more channels, and which are hopeless
until instrumentation changes---the same $\log\det$ and $\lambda_{\min}$ criteria
that drive D- and E-optimal design.

\paragraph{Stability as the master dial.}
The stability--identifiability law \eqref{eq:arlaw} explains several empirical
regularities of \TSFM s at once: strongly autocorrelated (near-nonstationary)
series are forecast with striking accuracy because their dominant modes carry
enormous Fisher information; white-ish series are nearly unforecastable because
they carry none; and the difficulty of a series is, to first order, its distance
from the unit circle. It also warns that the same near-nonstationarity that makes
a series learnable makes its FIM ill-conditioned, so second-order (natural-%
gradient / K-FAC) training and careful tokenization matter most exactly where
forecasting is most rewarding.

\paragraph{Tokenization revisited.}
\Cref{prop:quant} places \Chronos's ``regression via classification'' inside the
information geometry: quantization is data processing that can only lose Fisher
information, and the bin count $B$ trades vocabulary size against the information
retained about fine dynamics. The categorical head's advantage---arbitrary,
multi-modal predictive distributions---is orthogonal to this and remains; the
Fisher view simply prices the discretization.

\paragraph{Limitations.}
The exact spectral results are linear--Gaussian; nonlinear systems retain the
structure through the extended/ensemble tangent filter but lose the closed-form
$1/(1-r^2)$ law. Our \Chronos-lite is small and trained per stability level
rather than as a single cross-system foundation model; the cross-system,
zero-shot regime---where one model amortizes identification across many
systems---is the natural sequel, and \cref{prop:aisbound} suggests the shared
latent's AIS quality is what governs transfer. Finally, we analyzed the
one-step FIM; multi-step and rolling forecast risk propagate it through the
dynamics, amplifying the role of the slow modes further.

\section{Conclusion}
\label{sec:conclusion}
A time-series foundation model, a world-action model, and a
vision--language--action policy are three channels of one partially observed
process on a shared approximate-information-state latent, and their learning
problems are blocks of one Fisher information matrix. Deleting the action channel
collapses the two-block theory to the single system-identification block, which is
exactly what a forecaster optimizes; and for the linear--Gaussian instance that
block's spectrum is written by the system's stability margins and observability.
The slow modes near the unit circle carry the information and are learned first;
the fast modes carry little and are learned last; the unobserved modes carry none
and are never learned until a sensor is added. Stability is the master dial,
trading learnability against conditioning; observability is the ceiling. The same
Fisher matrix that measures how well an agent can learn its world measures how
well a foundation model can learn a time series---and its small eigenvalues, in
both cases, are the message.

\begin{thebibliography}{99}
\small\setlength{\itemsep}{1pt}
\bibitem{wamvla} \emph{World-Action Models and Vision--Language--Action Models:
  One Latent, One Fisher Information---Two Orthogonal Blocks}, companion
  manuscript (2025). Built on the approximate-information-state framework
  of~\cite{ais2022}.
\bibitem{ais2022} J.~Subramanian, A.~Sinha, R.~Seraj, and A.~Mahajan,
  \emph{Approximate information state for approximate planning and reinforcement
  learning in partially observed systems}, J.~Mach.~Learn.~Res., 23 (2022),
  pp.~1--83.
\bibitem{chronos} A.~F.~Ansari, L.~Stella, C.~Turkmen, X.~Zhang, P.~Mercado,
  H.~Shen, O.~Shchur, S.~S.~Rangapuram, S.~Pineda Arango, S.~Kapoor,
  J.~Zschiegner, D.~C.~Maddix, H.~Wang, M.~W.~Mahoney, K.~Torkkola,
  A.~G.~Wilson, M.~Bohlke-Schneider, and Y.~Wang, \emph{Chronos: learning the
  language of time series}, Trans.~Mach.~Learn.~Res.~(2024).
\bibitem{ljung1999} L.~Ljung, \emph{System Identification: Theory for the User},
  2nd ed., Prentice Hall, 1999.
\bibitem{cappe2005} O.~Capp\'e, E.~Moulines, and T.~Ryd\'en, \emph{Inference in
  Hidden Markov Models}, Springer, 2005.
\bibitem{phillips1987} P.~C.~B.~Phillips, \emph{Time series regression with a
  unit root}, Econometrica, 55 (1987), pp.~277--301.
\bibitem{transtrum2015} M.~K.~Transtrum, B.~B.~Machta, K.~S.~Brown,
  B.~C.~Daniels, C.~R.~Myers, and J.~P.~Sethna, \emph{Perspective: Sloppiness
  and emergent theories in physics, biology, and beyond}, J.~Chem.~Phys., 143
  (2015), 010901.
\bibitem{gutenkunst2007} R.~N.~Gutenkunst, J.~J.~Waterfall, F.~P.~Casey,
  K.~S.~Brown, C.~R.~Myers, and J.~P.~Sethna, \emph{Universally sloppy parameter
  sensitivities in systems biology models}, PLoS Comput.~Biol., 3 (2007), e189.
\bibitem{bellman1970} R.~Bellman and K.~J.~{\AA}str\"om, \emph{On structural
  identifiability}, Math.~Biosci., 7 (1970), pp.~329--339.
\bibitem{das2023} A.~Das, W.~Kong, R.~Sen, and Y.~Zhou, \emph{A decoder-only
  foundation model for time-series forecasting}, arXiv:2310.10688, 2023.
\bibitem{woo2024} G.~Woo, C.~Liu, A.~Kumar, C.~Xiong, S.~Savarese, and D.~Sahoo,
  \emph{Unified training of universal time series forecasting transformers},
  in Proc.~Int.~Conf.~Mach.~Learn.~(ICML), 2024.
\bibitem{rasul2023} K.~Rasul et al., \emph{Lag-Llama: towards foundation models
  for time series forecasting}, arXiv:2310.08278, 2023.
\bibitem{goswami2024} M.~Goswami, K.~Szafer, A.~Choudhry, Y.~Cai, S.~Li, and
  A.~Dubrawski, \emph{MOMENT: a family of open time-series foundation models},
  in Proc.~Int.~Conf.~Mach.~Learn.~(ICML), 2024.
\end{thebibliography}
\end{document}
